% SPDX-License-Identifier: GPL-3.0-or-later
% Copyright (C) 2026 Alessandro Greco
\chapter{Система каталогов}\label{cap:12-directory}
\aggiornatoal{23 сентября 2026 г.}{коммит torspec \texttt{8e1f245}; tor 0.4.9.13; метод консенсуса 36}

\begin{obiettivi}
К концу этой главы вы сможете:
\begin{itemize}[nosep]
  \item объяснять, зачем Tor нужна система каталогов и какие риски она должна предотвращать;
  \item различать её документы: дескрипторы, микродескрипторы, голоса, консенсусы и сертификаты центров каталогов;
  \item описывать, как центры каталогов каждый час голосуют и объединяют свои голоса;
  \item понимать критерии назначения флагов;
  \item читать консенсус построчно и проверять его подписи;
  \item загружать и проверять консенсус лабораторной сети с помощью собственного кода (лабораторная работа~6).
\end{itemize}
\end{obiettivi}

В предыдущих лабораторных работах мы всегда заранее знали, к каким реле подключаться, какие ключи они используют и какие порты слушают: эту информацию мы читали из файлов лабораторной сети. У реального клиента такой роскоши нет. Он должен узнать, какие реле существуют, где они находятся, какие ключи используют и насколько им можно доверять. Этим занимается \textbf{система каталогов}, которой посвящена эта глава.

% ---------------------------------------------------------------------------
\section{Проблема}\label{sec:dir-problema}

Система каталогов Tor должна удовлетворять нескольким частично конкурирующим требованиям:
\begin{description}
  \item[Корректность] Информация о реле (адреса и ключи) должна быть подлинной. Противник, убедивший клиента, что определённое реле имеет другой ключ, мог бы перехватить его трафик.
  \item[Единое представление] Все клиенты должны видеть \emph{один и тот же} список реле. Если бы противник мог выдавать каждому клиенту свой список, возможно содержащий только несколько контролируемых им реле, клиентов можно было бы узнавать по используемым ими реле. Это атака \emph{разделения (partitioning)}, также называемая \emph{эпистемической}.
  \item[Устойчивость к злоупотреблениям] Добавить реле в сеть может любой. Система должна ограничивать влияние того, кто добавляет множество реле (\emph{Sybil}-атака), и уметь исключать реле, ведущие себя некорректно.
  \item[Масштабируемость] Миллионы клиентов загружают каталожную информацию несколько раз в день. Её распространение должно быть экономичным.
\end{description}

Решение Tor, сформировавшееся за три версии протокола каталогов, --- небольшой набор доверенных серверов, \textbf{центры каталогов (directory authorities)}. Каждый час они голосуют и совместно подписывают единый документ: \textbf{консенсус}~\spec{dir-spec}{Outline}.

\begin{storia}[Три поколения]
В первой версии протокола каждый центр каталогов публиковал собственный список реле, поэтому нечестный центр каталогов мог вводить клиентов в заблуждение. Во второй клиенты загружали списки всех центров каталогов и сами объединяли представления, многократно увеличивая расход полосы. Третья версия, используемая более пятнадцати лет, передаёт объединение центрам каталогов: они создают один документ, подписанный всеми, а клиентам остаётся только проверить его. Тогда же центры каталогов перестали подписывать документы тем же ключом, который использовали как реле и который приходилось хранить в памяти в открытом виде.
\end{storia}

% ---------------------------------------------------------------------------
\section{Участники}\label{sec:dir-attori}

\begin{description}
  \item[Центр каталогов (Directory authority)] Серверы, которыми управляют доверенные участники сообщества Tor. В публичной сети их девять, плюс отдельный центр для мостов (\cref{cap:16-censura}). Их список встроен в tor (\file{src/app/config/auth\_dirs.inc}), поэтому каждый клиент начинает с одного и того же набора. Каждый центр имеет долгосрочный \emph{ключ идентичности}, хранящийся офлайн, который сертифицирует \emph{ключ подписи} со сроком действия от 3 до 12 месяцев. Последний подписывает голоса и консенсусы.
  \item[Кэш каталогов (Directory cache)] Реле, которые загружают консенсусы и дескрипторы у центров каталогов и повторно распространяют их клиентам, уменьшая нагрузку на центры каталогов. Они имеют флаг \texttt{V2Dir} и обслуживают потоки \file{BEGIN\_DIR} (\cref{cap:11-stream}).
  \item[Резервный сервер каталогов (Fallback directory)] Список примерно из двухсот стабильных реле, также встроенный в tor (\file{fallback\_dirs.inc}; 200 записей в tor 0.4.9.13). Новый клиент запрашивает первый консенсус у одного из них, а не перегружает центры каталогов.
  \item[Реле] Реле периодически публикуют собственные дескрипторы центрам каталогов.
  \item[Клиент] Клиенты загружают консенсус и микродескрипторы, проверяют их и используют для выбора маршрутов (\cref{cap:13-percorso}).
\end{description}

На \cref{fig:dir-flusso} показано движение этих документов.

\begin{figure}[htbp]
\centering
% SPDX-License-Identifier: GPL-3.0-or-later
% Copyright (C) 2026 Alessandro Greco
%% Русское альтернативное описание EPUB.
\altfig{Упрощённая схема каталогов Tor: ретрансляторы загружают дескрипторы известным центрам каталога; обновление обычно происходит примерно через 18 часов либо раньше при существенных изменениях. Стрелка к одному центру условна. Три значка обозначают центры каталога: они обмениваются голосами, вычисляют консенсус на основе доступных голосов и обмениваются подписями; при недостающих или разных голосах результаты могут расходиться. Каталожные кеши отдают консенсус микродескрипторов и сами микродескрипторы. Клиент проверяет подписи более половины доверенных центров и запрашивает недостающие микродескрипторы по хешу. Обычно новый консенсус публикуется каждый час.}
% Поток документов каталога.
\begin{tikzpicture}[font=\sffamily\small,>=Stealth,
  aut/.style={relay,fill=torviola!15,draw=torviola,minimum size=1cm,font=\sffamily\scriptsize},
  doc/.style={font=\sffamily\scriptsize,fill=white,inner sep=1pt,align=center}]
  \node[aut] (a1) at (0,2.2) {центр 1};
  \node[aut] (a2) at (1.8,3.2) {центр 2};
  \node[aut] (a3) at (3.6,2.2) {центр 3};
  \draw[<->,thick,torviola!70] (a1) -- (a2); \draw[<->,thick,torviola!70] (a2) -- (a3); \draw[<->,thick,torviola!70] (a1) -- (a3);
  \node[doc,text=torviola] at (1.8,1.75) {голоса и подписи};
  \node[relay] (r) at (-3.2,2.2) {ретранслятор};
  \draw[cellaf] (r) -- node[doc,above]{дескриптор\\(обычно 18 ч)} (a1);
  \node[relay,fill=torgrigio!10,draw=torgrigio] (c) at (1.8,-0.4) {\scriptsize кеш};
  \draw[risposta] (a1) -- node[doc,left=2pt]{консенсус,\\микродескрипторы} (c);
  \draw[risposta] (a3) -- (c);
  \node[client] (cl) at (6.2,-0.4) {Клиент};
  \draw[risposta] (c) -- node[doc,above]{консенсус microdesc\\(примерно раз в час)} (cl);
  \draw[risposta] (c) to[bend right=25] node[doc,below]{недостающие микродескрипторы, по хешу} (cl);
  \node[font=\sffamily\scriptsize,text=torgrigio,align=center] at (6.2,2.6) {проверка: подписи\\более половины доверенных центров};
  \draw[dotted,torgrigio] (6.2,2.2) -- (cl);
\end{tikzpicture}

\caption{Поток документов системы каталогов. Реле публикуют дескрипторы центрам каталогов; центры каталогов обмениваются голосами, вычисляют консенсус на основе полученных голосов и обмениваются подписями; кэши каталогов загружают консенсусы и микродескрипторы и обслуживают ими клиентов.}
\label{fig:dir-flusso}
\end{figure}

% ---------------------------------------------------------------------------
\section{Документы}\label{sec:dir-documenti}

Все документы каталогов представляют собой текст в формате, который спецификация называет \emph{netdoc}: последовательность строк (\emph{элементов (items)}), каждая с ключевым словом и аргументами, иногда с последующим объектом в base64 между строками \texttt{-----BEGIN ...-----} и \texttt{-----END ...-----}. Этот формат легко читать человеку и разбирать программой, что мы увидим в лабораторной работе.

\begin{description}
  \item[Дескриптор реле] (\emph{server descriptor}) Создаётся и подписывается реле: адреса и порты, ключи (идентичность, ключ onion ntor, сертификаты Ed25519), заявленная пропускная способность, полная политика выхода, семейство, платформа и версия. Он повторно публикуется не реже чем каждые 18 часов (\file{FORCE\_REGENERATE\_DESCRIPTOR\_INTERVAL} в коде tor) или раньше при существенных изменениях.
  \item[Документ extra-info] Статистика, не нужная для работы протокола, например трафик мостов по странам, публикуется отдельно, чтобы не раздувать дескрипторы. Эти данные питают Tor Metrics (\cref{cap:24-metrics}).
  \item[Микродескриптор] Подмножество дескриптора, действительно нужное клиенту и редко меняющееся: ключ onion ntor, семейство, сводка политики выхода и идентичность Ed25519. Его вычисляют центры каталогов; каждый микродескриптор идентифицируется дайджестом SHA-256 своего текста.
  \item[Голос] Оценка сети одного центра каталогов: флаги, пропускная способность и версия каждого реле.
  \item[Консенсус] Результат объединения голосов, подписанный центрами каталогов. Он имеет два \emph{варианта (flavor)}: \texttt{ns}, содержащий полный дайджест дескриптора для каждого реле, и \texttt{microdesc}, содержащий дайджест микродескриптора. Современные клиенты используют только более компактный второй вариант.
  \item[Сертификат ключа центра каталогов] Связывает ключ подписи центра каталогов с его идентичностью.
\end{description}

\begin{approfondimento}[Зачем микродескрипторы?]
Полный дескриптор имеет размер несколько килобайт и часто меняется, например при изменении заявленной полосы. Если бы клиент несколько раз в день загружал все дескрипторы тысяч реле, это потребовало бы значительной полосы, особенно на мобильных соединениях. Микродескриптор содержит лишь редко меняющуюся информацию. Часто меняющиеся сведения, такие как пропускная способность и флаги, находятся в консенсусе, который клиент всё равно загружает. Поскольку клиент запрашивает микродескрипторы по дайджесту, а не по идентичности реле, кэш каталогов не может выдать одному клиенту иной микродескриптор, чем всем остальным.
\end{approfondimento}

% ---------------------------------------------------------------------------
\section{Голосование}\label{sec:dir-voto}

\subsection{Расписание}

У каждого консенсуса три временные метки~\spec{dir-spec}{Voting timeline}:
\begin{description}
  \item[\texttt{valid-after}] момент, когда он становится действительным;
  \item[\texttt{fresh-until}] момент, когда он перестаёт быть самым свежим; обычно совпадает с \texttt{valid-after} следующего консенсуса;
  \item[\texttt{valid-until}] момент, после которого клиент больше не может его использовать.
\end{description}
В публичной сети интервал голосования равен одному часу, а консенсус действует три интервала (опции \opt{V3AuthVotingInterval} и \opt{V3AuthNIntervalsValid}). Поэтому в любой момент одновременно существует несколько действительных консенсусов, и задержка центрам каталогов не блокирует клиентов. До начала каждого интервала центры каталогов имеют пять минут на обмен голосами и пять минут на обмен подписями (\opt{V3AuthVoteDelay} и \opt{V3AuthDistDelay}). В лабораторной сети всё ускорено: голосование каждые 20 секунд, по 4 секунды на каждую фазу (\cref{sec:rp-opzioni}).

\subsection{Объединение голосов}

Каждый центр каталогов собирает голоса остальных и вычисляет консенсус детерминированным алгоритмом, чтобы все центры каталогов получили одинаковый документ и могли его подписать~\spec{dir-spec}{Computing a consensus from a set of votes}. Основные правила:
\begin{itemize}
  \item реле появляется в консенсусе только если его знает достаточное число центров каталогов и, в современных методах консенсуса (24 и новее; лаборатория использует 36), только если оно имеет одновременно \texttt{Running} и \texttt{Valid};
  \item реле получает флаг, если его назначает более половины центров, голосующих по этому флагу;
  \item пропускная способность (\texttt{w}) --- нижняя медиана проголосованных значений, кроме случая, когда как минимум три центра каталогов \emph{измерили} пропускную способность: тогда используется медиана этих измерений;
  \item сетевые параметры (\texttt{params}) включаются, если за них голосует большинство всех центров каталогов либо как минимум три центра каталогов (\texttt{i > total\_authorities/2 || i >= 3} в коде); выбирается нижняя медиана;
  \item временные метки --- медианы предложенных значений.
\end{itemize}
Поскольку центры каталогов могут работать на разных версиях tor, документ объявляет используемый \emph{метод консенсуса}: число, идентифицирующее точный набор правил, выбранное как наивысший метод, поддерживаемый достаточным числом центров каталогов. Консенсус лабораторной сети использует метод 36.

\subsection{Флаги}\label{sec:dir-flag}

Флаги кратко выражают оценку реле каждого центра каталогов~\spec{dir-spec}{Assigning flags in a vote}. Спецификация предупреждает, что точные критерии могут меняться и клиенты не должны от них зависеть. \Cref{tab:flag} сводит современные критерии.

\begin{table}[htbp]
\centering\small
\begin{tabularx}{\textwidth}{@{}l>{\raggedright\arraybackslash}X@{}}
\toprule
Флаг & Критерий (кратко) \\
\midrule
\texttt{Running} & центр каталогов установил соединение с реле на всех его ORPorts в последние 45 минут \\
\texttt{Valid} & версия tor не известна как дефектная; реле не исключено центрами каталогов \\
\texttt{Fast} & пропускная способность находится среди самых быстрых семи восьмых либо не ниже 100\,KB/s \\
\texttt{Stable} & взвешенное по времени среднее время между отказами не ниже медианы либо не меньше 7 дней \\
\texttt{Guard} & \texttt{Fast}, \texttt{Stable}, \emph{familiar} (в сети несколько недель), хорошая доступность, достаточная пропускная способность, \texttt{V2Dir} \\
\texttt{Exit} & разрешает выходной трафик как минимум к одному /8-блоку одновременно на портах 80 и 443 \\
\texttt{HSDir} & \texttt{Stable} и \texttt{Fast}, работает не менее 96 часов: может хранить дескрипторы onion-сервисов \\
\texttt{V2Dir} & обслуживает каталожную информацию (кэш каталогов) \\
\texttt{Authority} & является центром каталогов \\
\texttt{BadExit} & выходное реле, которое не следует использовать для выхода, например из-за изменения трафика \\
\texttt{MiddleOnly} & следует использовать только как промежуточное реле (введён в 0.4.7.2-alpha) \\
\texttt{StaleDesc} & дескриптор опубликован более 18 часов назад \\
\texttt{Sybil} & превышает максимальное число реле на один IP-адрес (обычно 2) \\
\bottomrule
\end{tabularx}
\caption{Флаги консенсуса и краткие критерии их назначения центрами каталогов.}
\label{tab:flag}
\end{table}

\texttt{BadExit} и \texttt{MiddleOnly} позволяют центрам каталогов реагировать на подозрительные реле, не удаляя их полностью. Флаг \texttt{Sybil} ограничивает влияние того, кто запускает множество реле с одного адреса. В лабораторной сети, где все реле работают на 127.0.0.1, ограничение отключено параметром \texttt{AuthDirMaxServersPerAddr=0}; это будет видно в строке \texttt{params}.

\subsection{Пропускная способность}

Какова пропускная способность реле? Значению, заявленному в его дескрипторе, доверять нельзя: вредоносное реле могло бы объявить огромную пропускную способность, чтобы его выбирали чаще. Некоторые центры каталогов, называемые \emph{центрами измерения пропускной способности (bandwidth authorities)}, поэтому измеряют пропускную способность реле, пропуская через них трафик специальным инструментом (сейчас \emph{sbws}, \cref{cap:24-metrics}), и включают результат в голоса как \texttt{Measured}. Если доступно не менее трёх измерений, консенсус использует их медиану. Иначе строка \texttt{w} содержит \texttt{Unmeasured=1} и ограниченное значение. Consensus bandwidth --- не просто скорость, а \emph{вес}: он определяет вероятность выбора реле (\cref{cap:13-percorso}).

% ---------------------------------------------------------------------------
\section{Анатомия консенсуса}\label{sec:dir-anatomia}

Ниже приведён сокращённый консенсус варианта microdesc лабораторной сети, загруженный в лабораторной работе~6. Фрагменты взяты из сохранённой в репозитории копии (\file{labs/ru/код/torlab/test/data/consensus-microdesc.txt}), используемой тестами лаборатории. Поэтому их можно изучать построчно без запуска сети. Копия была получена 23 сентября 2026 года, когда лабораторные реле ещё работали на tor 0.4.9.12. Именно эта версия находится в строке \texttt{v} записи реле, а не референсная версия книги, на которой сеть работает сегодня.

\paragraph{Заголовок.}
\begin{lstlisting}[numbers=none]
network-status-version 3 microdesc
vote-status consensus
consensus-method 36
valid-after 2026-09-23 15:37:20
fresh-until 2026-09-23 15:37:40
valid-until 2026-09-23 15:38:00
voting-delay 4 4
known-flags Authority Exit Fast Guard HSDir ... Running Stable ... Valid
params AuthDirMaxServersPerAddr=0 cc_alg=2
\end{lstlisting}
Строка \texttt{params} содержит \emph{параметры консенсуса}, с помощью которых центры каталогов меняют поведение всей сети без обновления программного обеспечения. Здесь, помимо ограничения на один адрес, задаётся алгоритм управления перегрузкой (\file{cc\_alg=2}, Vegas, \cref{sec:cc}). В публичной сети параметров несколько десятков.

\paragraph{Запись реле.}
\begin{lstlisting}[numbers=none]
r test006e ml1aPg9UFy4VCx85YcpsRs/9shc 2038-01-01 00:00:00 127.0.0.1 5106 0
m FfrgO+s4VYS6WKBq3pjfMDs+uZ2uVsvgG88JilbhIOo
s Exit Fast Guard HSDir Running Stable V2Dir Valid
v Tor 0.4.9.12
pr Conflux=1 Cons=1-2 Desc=1-4 ... Link=3-5 ... Relay=2-6
w Bandwidth=4096
\end{lstlisting}
\begin{description}
  \item[\texttt{r}] псевдоним реле, дайджест RSA-идентичности (base64), поле даты, больше не имеющее смысла в консенсусе, IPv4-адрес, ORPort и DirPort;
  \item[\texttt{m}] дайджест микродескриптора, который следует загрузить;
  \item[\texttt{s}] флаги;
  \item[\texttt{v}, \texttt{pr}] версия tor и поддерживаемые \emph{подпротоколы}: здесь клиент, например, узнаёт, что реле поддерживает CGO (\texttt{Relay=6}) или версии протокола соединения 3--5;
  \item[\texttt{w}] вес пропускной способности в килобайтах в секунду.
\end{description}
В варианте \texttt{ns} каждая запись также содержит строку \texttt{p}, кратко описывающую политику выхода по портам. В варианте microdesc это описание находится в микродескрипторе и не повторяется в консенсусе.

\paragraph{Нижняя часть и подписи.}
\begin{lstlisting}[numbers=none]
directory-footer
bandwidth-weights Wbd=821 Wbe=0 Wbg=5000 Wbm=10000 Wdb=10000 ...
directory-signature sha256 3776EF73A9214587AEFB3D732FA4A443CFF491BA
                           79E102D226590BC5D1D1EBBAD8A16122EC0F9801
-----BEGIN SIGNATURE-----
...
-----END SIGNATURE-----
\end{lstlisting}
Значения \texttt{bandwidth-weights} --- коэффициенты, которыми клиенты балансируют использование реле в разных позициях цепи (\cref{cap:13-percorso}). Далее следуют подписи, по одной от каждого центра каталогов: алгоритм хеширования, дайджест идентичности центра каталогов, дайджест его ключа подписи и подпись.

% ---------------------------------------------------------------------------
\section{Проверка консенсуса}\label{sec:dir-firme}

Клиент принимает консенсус только при наличии действительных подписей \textbf{более чем половины} известных ему центров каталогов. В коде tor порог равен \file{n\_v3\_authorities/2 + 1} (\file{src/feature/nodelist/networkstatus.c}): 5 подписей из 9 в публичной сети или 2 из 3 в лаборатории. Для проверки подписи клиенту нужны:
\begin{enumerate}
  \item \textbf{сертификат} центра каталогов, который клиент загружает (\nolinkurl{/tor/keys/all}) и проверяет по заранее известной идентичности центра каталогов;
  \item \textbf{подписанный текст}: документ от строки \texttt{network-status-version} до пробела включительно после первого ключевого слова \texttt{directory-signature};
  \item схема подписи: RSA с дополнением PKCS\#1 v1.5, применённым непосредственно к дайджесту (SHA-1 или SHA-256), \emph{без} идентификатора алгоритма, обычно ожидаемого PKCS\#1~\spec{dir-spec}{Signing documents}.
\end{enumerate}

\begin{approfondimento}[До пробела или до перевода строки?]
Общий раздел спецификации о подписанных документах говорит, что подпись охватывает текст до \emph{перевода строки} после строки подписи. Раздел, относящийся именно к консенсусу, вместо этого указывает \emph{пробел} после \texttt{directory-signature}, чтобы все центры каталогов подписывали одинаковые байты, хотя строки их подписей различаются. В лабораторной работе~6 мы проверили обе трактовки на реальном консенсусе лабораторной сети. Все три подписи были действительны при правиле пробела; ни одна --- при правиле перевода строки. При сомнении экспериментальная проверка против независимой реализации остаётся окончательным арбитром.
\end{approfondimento}

\begin{attenzione}[Список центров каталогов --- корень доверия]
Вся система опирается на список центров каталогов, встроенный в программное обеспечение. Пользователь может его изменить, но спецификация не рекомендует этого делать: клиент с другим набором центров каталогов увидит отличный от остальных консенсус и станет различимым. Сами центры каталогов являются целями. Если большинство будет скомпрометировано, они смогут подписать консенсус с вредоносными реле. Поэтому центров каталогов управляются разными людьми в разных странах, а ключи идентичности хранятся офлайн.
\end{attenzione}

% ---------------------------------------------------------------------------
\section{Как клиент остаётся актуальным}\label{sec:dir-client}

Клиент не загружает новый консенсус ровно каждый час: если бы все клиенты делали это одновременно, кэши каталогов оказались бы перегружены. Вместо этого он выбирает случайный момент между тремя четвертями первого интервала после потери консенсусом свежести и семью восьмыми оставшегося времени до истечения~\spec{dir-spec}{Client operation}. Спецификация приводит пример: для консенсуса, действующего с 1:00, свежего до 2:00 и действительного до 4:00, клиент начнёт искать новый в случайный момент между 2:45 и 3:50.

Для экономии полосы пропускания клиент может запросить только \emph{различия} с уже имеющимся консенсусом (\emph{consensus diff}), идентифицируя его дайджестом SHA3-256. Затем он загружает только отсутствующие микродескрипторы, запрашивая их по дайджесту (\nolinkurl{/tor/micro/d/}, после чего дайджесты разделяются \texttt{-}).

% ---------------------------------------------------------------------------
\section{Лабораторная работа: загрузка и проверка консенсуса}\label{sec:lab-directory}

\begin{laboratorio}[6]{Консенсус лабораторной сети}
\textbf{Цель.} Написать \file{labs/ru/код/torlab/directory.py} для загрузки консенсуса, сертификатов и микродескрипторов у центра каталогов через цепь, их разбора и проверки подписей.\\
\textbf{Предварительные условия.} Лабораторные работы~1--5; работающая лабораторная сеть.\\
\textbf{Оценочное время.} 3 часа.
\end{laboratorio}

\paragraph{Шаг 1: загрузка через BEGIN\_DIR.} Как tor, используем одношаговую цепь к серверу каталогов без управления перегрузкой и открываем потоки \file{BEGIN\_DIR}. Внутри потока используется HTTP/1.0:
\begin{lstlisting}[language=Python]
def download(circ, path, stream_id):
    s = Stream(circ, stream_id)
    circ.send(RelayMessage(RelayCommand.BEGIN_DIR, stream_id, b""))
    _, msg = s._next()                      # CONNECTED, пустое тело
    s.send(f"GET {path} HTTP/1.0\r\n\r\n".encode())
    header, _, body = s.receive_all().partition(b"\r\n\r\n")
    return body

circ = build([relay("000a")], cc=False)
text = download(circ, "/tor/status-vote/current/consensus-microdesc", 1)
\end{lstlisting}

\paragraph{Шаг 2: формат netdoc.} Функция примерно из двадцати строк, \texttt{items}, делит документ на элементы и собирает объекты \texttt{-----BEGIN}/\texttt{-----END}. На её основе строятся два парсера: \file{parse\_consensus}, создающий запись каждого реле с флагами, пропускной способностью, политикой выхода и дайджестом микродескриптора, и \file{parse\_certificates}.

\paragraph{Шаг 3: проверка подписей.}
\begin{lstlisting}[language=Python]
def signed_text(text):
    start = text.index("network-status-version ")
    pos = text.index("\ndirectory-signature ") + len("\ndirectory-signature ")
    return text[start:pos].encode()

def verify_signature(consensus, signature, cert):
    h = hashlib.sha256 if signature.algorithm == "sha256" else hashlib.sha1
    expected = h(signed_text(consensus.text)).digest()
    recovered = cert.signing_key.recover_data_from_signature(
        signature.signature, padding.PKCS1v15(), None)
    return recovered == expected
\end{lstlisting}
Метод \file{recover\_data\_from\_signature} с алгоритмом \texttt{None} возвращает дайджест, встроенный в подпись, не ожидая идентификатора алгоритма: именно так устроена схема Tor.

\paragraph{Шаг 4: результаты.} В одном из наших запусков в лабораторной сети (даты, имена и количества зависят от экземпляра сети):
\begin{lstlisting}[numbers=none]
консенсус: 6035 байт; 2026-09-23 15:37:20 -> 2026-09-23 15:38:00
метод 36; подписи 3; сертификаты 3
params: {'AuthDirMaxServersPerAddr': 0, 'cc_alg': 2}
  test006e 127.0.0.1:5106 w=4096 flags=['Exit', 'Fast', 'Guard', ...]
  test000a 127.0.0.1:5100 w=20   flags=['Authority', 'Fast', 'Guard', ...]
  ...
действительных подписей: 3 (необходимо: 2)
микродескрипторов: 8; соответствуют строкам m: True
\end{lstlisting}
В исходной записи лабораторного опыта консенсус охватывает 40 секунд, то есть два 20-секундных интервала, подписан всеми тремя центрами каталогов, а дайджест SHA-256 каждого загруженного микродескриптора совпадает со строкой \texttt{m} в консенсусе. Цепочка доверия проходит от ключей центров каталогов до ключа ntor каждого реле.

\paragraph{Шаг 5: отрицательные проверки.} Проверка, которая никогда не завершается ошибкой, ничего не проверяет. Измените одну цифру консенсуса (\texttt{Bandwidth=4096} на \texttt{4097}), и число действительных подписей упадёт до нуля; при наличии только одного известного центра каталогов --- до единицы. Файл \file{test/test\_directory.py} повторяет эти проверки по сохранённой копии консенсуса, поэтому тесты работают и без сети.

\paragraph{Шаг 6: настоящий консенсус.} Тот же код работает с публичной сетью. Загрузите через CollecTor консенсус и соответствующие сертификаты центров каталогов за тот же период и проверьте:
\begin{lstlisting}[language=Python,numbers=none]
c = parse_consensus(open("consensus-microdesc-2026-09-23-15").read())
cert = parse_certificates(open("certs").read())
valid, required, _ = verify_consensus(c, cert, 9)
print(valid, required)          # 9 5
\end{lstlisting}
В приведённом примере действительны все девять подписей, что выше порога в пять. При первой попытке мешает одна деталь: файлы CollecTor начинаются строкой-аннотацией \texttt{@type network-status-microdesc-consensus-3 1.0}, которая не входит в подписанный документ. Если включить её, дайджест изменится и \emph{все} подписи станут недействительными. Это выглядит как криптографическая ошибка, но на деле является ошибкой форматирования. Поэтому функция \texttt{signed\_text} начинает с \texttt{network-status-version}, и отдельный тест проверяет это поведение.

% ---------------------------------------------------------------------------
\section*{Итоги}
\begin{itemize}
  \item Система каталогов должна давать всем клиентам одинаковое подлинное представление сети, противостоять Sybil-атакам и масштабироваться до миллионов пользователей.
  \item Девять центров каталогов плюс один для мостов, перечисленные в коде, голосуют каждый час и подписывают единый консенсус, действительный три часа.
  \item Реле публикуют дескрипторы; клиенты используют консенсус варианта microdesc и микродескрипторы, загружаемые по дайджесту.
  \item Флаг назначается, если за него голосует более половины центров каталогов, ответственных за этот флаг; пропускная способность, измеренная как минимум тремя центрами каталогов, имеет приоритет над заявленной.
  \item Клиент принимает консенсус с действительными подписями более половины центров каталогов; подпись охватывает текст до пробела включительно после \texttt{directory-signature}.
  \item Параметры консенсуса изменяют поведение всей сети без обновления программного обеспечения.
\end{itemize}

% ---------------------------------------------------------------------------
\section*{Упражнения}

\begin{esercizio}[1]\label{es:12-validita}
Консенсус имеет \texttt{valid-after} 10:00, \texttt{fresh-until} 11:00 и \texttt{valid-until} 13:00. В каком интервале клиент будет искать замену? Что произойдёт, если к 13:00 нового консенсуса не будет?
\end{esercizio}

\begin{esercizio}[1]\label{es:12-flag}
Реле имеет пропускную способность 500\,KB/s, находится в сети два дня и использует политику выхода \texttt{accept *:443, reject *:*}. Какие флаги оно может получить? Почему не \texttt{Exit}?
\end{esercizio}

\begin{esercizio}[2]\label{es:12-partizione}
Опишите атаку разделения (partitioning), в которой вредоносный кэш каталогов отправляет конкретному клиенту старый, но ещё действительный консенсус. Что ограничивает атаку? Почему запрос микродескрипторов по дайджесту помогает?
\end{esercizio}

\begin{esercizio}[2]\label{es:12-mediana}
Три центра каталогов измеряют пропускную способность реле: 5000, 5200 и 90000. Какое значение попадёт в консенсус? Что изменится, если третий центр каталогов, вступивший в сговор с реле, будет единственной, кто его измерял?
\end{esercizio}

\begin{esercizio}[2]\label{es:12-voti}
Найдите файл \file{v3-status-votes} в каталоге одного центра каталогов лабораторной сети. Сравните её голос с консенсусом: какие строки отличаются и почему?
\end{esercizio}

\begin{esercizio}[2]\label{es:12-md}
Выберите реле из консенсуса, загруженного в лабораторной работе. Вручную вычислите дайджест его микродескриптора с помощью \cmd{sha256sum} и сравните со строкой \texttt{m}. С какой строки начинается хешируемый текст и где заканчивается?
\end{esercizio}

\begin{esercizio}[3]\label{es:12-autorita}
Остановите один из трёх центров каталогов лабораторной сети и через минуту снова загрузите консенсус. Сколько подписей он содержит? Принимает ли его ваша проверка? Повторите с двумя остановленными центрами каталогов и объясните наблюдаемое после истечения консенсуса.
\end{esercizio}

\begin{esercizio}[3]\label{es:12-ns}
Также загрузите консенсус \texttt{ns} (\nolinkurl{/tor/status-vote/current/consensus}) и проверьте его подписи. Какие алгоритмы хеширования он использует? Какая дополнительная информация находится в строке \texttt{r} по сравнению с вариантом microdesc?
\end{esercizio}

% ---------------------------------------------------------------------------
\section*{Дополнительная литература}
\begin{itemize}
  \item \titolo{Tor directory protocol, version 3} (dir-spec)~\cite{torspec}: форматы, голосование, назначение флагов и работа клиентов и кэшей каталогов.
  \item Документ протокола общей случайности (shared randomness) (srv-spec), с помощью которого центры каталогов ежедневно создают общее случайное значение; оно понадобится для onion-сервисов (\cref{cap:14-onion-services}).
  \item Tor Metrics и CollecTor, \url{https://metrics.torproject.org/}: исторические архивы консенсусов, голосов и дескрипторов.
\end{itemize}
